\addcontentsline{toc}{chapter}{Wykaz skrótów}
\section*{Wykaz skrótów}

\begin{longtable}{cp{0.8\textwidth}}

AI    & Artificial Intelligence - sztuczna inteligencja. \\
API   & Application Programming Interface - interfejs programowania aplikacji. \\
FSM   & Finite State Machine - skończona maszyna stanów. \\
HUD   & Heads-Up Display - wyświetlacz przeziernikowy; elementy UI widoczne podczas rozgrywki. \\
MIT   & Massachusetts Institute of Technology (licencja open-source MIT). \\
NPC   & Non-Player Character - postać niezależna (niegrywalny bohater). \\
PC    & Personal Computer - komputer osobisty. \\
RPG   & Role-Playing Game - gra fabularna. \\
UI    & User Interface - interfejs użytkownika. \\
UX    & User Experience - doświadczenie użytkownika. \\
VR    & Virtual Reality - wirtualna rzeczywistość. \\
AR    & Augmented Reality - rozszerzona rzeczywistość. \\
2D    & Two-Dimensional - dwuwymiarowy. \\
3D    & Three-Dimensional - trójwymiarowy. \\
2.5D  & Styl graficzny łączący elementy 2D i 3D; środowisko 3D z rozgrywką na płaszczyźnie 2D. \\
AAA   & Triple-A - produkcja gry z dużym budżetem przez duże studio. \\
Indie & Independent - niezależna produkcja gry (bez wydawcy). \\
HP    & Hit Points / Health Points - punkty życia postaci. \\
EXP   & Experience Points - punkty doświadczenia. \\
OOP   & Object-Oriented Programming - programowanie obiektowe. \\
.NET  & Platforma programistyczna firmy Microsoft (używana z C\#). \\
JSON  & JavaScript Object Notation - tekstowy format wymiany i zapisu danych. \\
DTO   & Data Transfer Object - obiekt transferu danych do serializacji/deserializacji. \\
PCG   & Procedural Content Generation - proceduralna generacja zawartości gry. \\
DDA   & Dynamic Difficulty Adjustment - dynamiczne dostosowanie trudności. \\
RL    & Reinforcement Learning - uczenie wzmocnione. \\
MAB   & Multi-Armed Bandit - problem bandyty wieloramiennego. \\
EMA   & Exponential Moving Average - wykładniczo ważona średnia ruchoma. \\
DPS   & Damage Per Second - obrażenia na sekundę. \\
EQ    & Equipment - ekwipunek. \\
MVP   & Minimum Viable Product - minimalny produkt akceptowalny. \\
SDK   & Software Development Kit. \\
LLM   & Large Language Model - duży model językowy (asystent AI w pracy nad projektem). \\
AABB  & Axis-Aligned Bounding Box - prostopadłościan ograniczający wyrównany do osi, używany w sortowaniu głębi sprite'ów. \\
UV    & Texture coordinates - współrzędne tekstury, używane w shaderach \texttt{character\_hurt} i \texttt{minimap\_mask}. \\
A*    & A-star - algorytm wyznaczania ścieżki na grafie, w grze realizowany przez \texttt{NavigationAgent3D}. \\
MDP   & Markov Decision Process - proces decyzyjny Markowa; pełny formalizm RL, do którego bandyta kontekstowy jest uproszczeniem. \\

\end{longtable}
